\section{Introduction}
\label{sec:intro}

The rapid proliferation of specialized, pre-trained neural networks has driven significant interest in methods that combine their distinct capabilities into a unified system~\cite{wortsman2022model, choshen2022fusing}. A prominent paradigm to achieve this is \textit{model merging}, where the parameters of independently fine-tuned specialist models (often low-rank adapters like LoRA~\cite{hu2021lora}) are blended or routed dynamically depending on the input sample~\cite{ilharco2022editing, yadav2023ties}. However, despite the potential of ensembling, the literature has experienced a rapid expansion of algorithmic complexity. Recent state-of-the-art dynamic ensembling methods rely on increasingly complex, over-parameterized neural networks—ranging from wave-superposition routing layers~\cite{gong2024qmerge} to multi-layer classifiers—to compute dynamic ensembling coefficients.

These parametric ensembling routers introduce severe systems-level and statistical bottlenecks. Defending these complex architectures requires introducing: (i) specialized calibration splits (e.g., small datasets to train the ensembling layers offline); (ii) multi-epoch AdamW optimization loops~\cite{kingma2014adam}; and (iii) hyperparameter heuristics to govern weight decay and learning rate schedules. This trend of adding layers of parameterization and training routines runs counter to the fundamental goal of ensembling: to leverage already pre-trained and optimized specialists with minimal overhead.

In this work, we take a stand against this unnecessary complexity, guided by Occam's razor. We ask a fundamental question: \textit{Can we perform high-fidelity, sample-wise dynamic model ensembling without training a single parameter or using a single calibration sample?} 

To answer this in the affirmative, we present and analyze \textbf{Parameter-Free Task-Space Projection (PFSR)}, a remarkably simple, elegant, and training-free method designed to perform dynamic model merging in a single online forward pass. Instead of training a routing model, PFSR extracts representative "task centroids" directly from the classification heads of pre-trained expert models, and projects online representations onto these centroids to guide dynamic ensembling via a temperature-scaled Softmax gating function.

To investigate whether we can further improve PFSR by decoupling task coordinates under task overlap, we propose and analyze an advanced extension: \textbf{L{\"o}wdin-Orthogonalized Task-Space Projection (OTSP)}. OTSP applies \textit{L{\"o}wdin Symmetric Orthogonalization}~\cite{lowdin1950non} to the task centroids offline. This operation computes the symmetric inverse square root of the Gram overlap matrix, producing an orthonormal task basis that is mathematically closest to the original directions in a least-squares sense, while preserving perfect order-invariance and symmetry across experts. At runtime, input features are projected onto this orthonormal basis to produce decoupled task coordinates.

Our empirical and theoretical investigation, conducted across a 10-seed rigorous simulation sandbox, yields profound and self-critical insights into ensembling dynamics. First, we show that under symmetric task correlations, OTSP and PFSR make \textit{exactly identical routing decisions}, proving that Löwdin orthogonalization is mathematically redundant. Second, under asymmetric task overlaps, we demonstrate that OTSP systematically \textit{underperforms} PFSR by 0.2\% to 1.6\% under noise. This is due to the \textit{Noise Amplification Penalty} and \textit{Noise Spillover Penalty}: the near-singular overlap makes $S$ ill-conditioned, scaling up the online projection coordinate noise variance and spilling noise from corrupted specialists onto clean coordinate axes. Consequently, the simpler PFSR is not only computationally cheaper, but systematically more robust than the more complex OTSP.

Third, we analyze a major numerical instability that can occur under sample-wise vectorized online deployment ($B=1$), which we call \textbf{Vectorization Collapse}. We show that when processed one sample at a time, unregularized linear routers collapse to 55.57\% accuracy due to the lack of constraint normalization and small-sample inductive overfitting on the small calibration split. We explicitly highlight that this instability is a consequence of omitting a probability-simplex constraint, and that any standard router employing simplex normalization (e.g., Softmax) is naturally immune to this issue.

Under a perfectly disjoint orthogonal subspace, PFSR achieves 74.46\% joint classification accuracy. While this matches the expert ceiling reference exactly, we transparently deconstruct this result under the \textbf{Orthogonal Masking Effect}: in disjoint sandboxes, any positive weight assigned to the correct expert (including simple Uniform Merging) yields the exact same classification prediction, meaning classification accuracy is flat and uninformative. We show instead that the true value of PFSR and OTSP lies in their high routing accuracy ($\ge 98.8\%$) and zero-parameter nature. Furthermore, we show that initializing the weights of a parametric Softmax routing network to exactly zero acts as an implicit maximum-entropy prior that shields it from small-sample overfitting.

Our core contributions are as follows:
\begin{itemize}
    \item We propose \textbf{Parameter-Free Task-Space Projection (PFSR)}, a 100\% training-free, data-free, closed-form ensembling router that extracts task centroids to perform dynamic model merging.
    \item We analyze \textbf{L{\"o}wdin-Orthogonalized Task-Space Projection (OTSP)} as an advanced extension to decouple task axes, proving mathematically and empirically that OTSP is redundant to PFSR under symmetric layouts and underperforms PFSR under asymmetric overlaps due to the Noise Amplification and Noise Spillover Penalties.
    \item We deconstruct the numerical instability of unconstrained linear routers under sample-wise vectorized ($B=1$) streaming, and reframe it as a pedagogical demonstration of the mathematical necessity of the probability-simplex constraint.
    \item We transparently analyze the \textbf{Orthogonal Masking Effect} in disjoint orthogonal sandboxes, demonstrating why joint classification accuracy is flat across all simplex-constrained methods and establishing routing accuracy as the primary informative metric for evaluation.
    \item We provide a mathematical and empirical analysis showing that simple zero-initialization of Softmax routing parameters acts as an implicit, perfect regularizer, challenging the need for complex, over-parameterized gating networks.
\end{itemize}
